% =====================================================================
\chapter{Problem Formulation}
\label{ch:problem}
% =====================================================================

This chapter turns the gaps identified in the literature review into a
precise problem. It describes the problem domain and the stakeholders,
states the problem in plain terms and in formal terms, depicts the inputs,
processes and outputs of the intended system, and lists the objectives
together with measurable success criteria.

% ---------------------------------------------------------------------
\section{Description of the Problem Domain}
% ---------------------------------------------------------------------

\subsection{Stakeholders and Their Requirements}

The intended users of an AI-text detector have different but overlapping
needs.
\begin{itemize}
  \item \textbf{Faculty and examiners} need to screen a whole batch of
        submissions at the end of an assignment cycle, quickly and without
        sending student work to an external service. When a submission is
        flagged they need to understand why, so that they can discuss it with
        the student.
  \item \textbf{Students} need protection from false accusations. A detector
        that flags honest work, and especially one that does so more often for
        some groups of students than for others, causes real harm.
  \item \textbf{Institutions} need a tool whose cost does not grow with the
        number of pages checked, which can run on ordinary hardware, and whose
        decisions can be audited if a case is appealed.
  \item \textbf{Publishers and platforms} need throughput: a detector that can
        triage large volumes of text and pass only questionable items to
        human moderators.
\end{itemize}
Taken together, these requirements describe a detector that is
\emph{accurate}, \emph{cheap per document}, \emph{explainable},
\emph{robust to new generators} and \emph{fair}. Chapter~\ref{ch:lit} showed
that no existing system offers all five properties.

\subsection{Identified Research Gaps}

The literature review leads to five interlinked limitations of current
detection tools. Each is stated below as a problem (P1--P5).

\textbf{P1 -- The black-box problem.} Leading commercial platforms such as
Turnitin and Copyleaks return a probability or a percentage without
disclosing the features or the decision process behind it. Academic
institutions cannot audit individual decisions, and a student has no concrete
evidence to respond to.

\textbf{P2 -- Cost and licensing barriers.} Institutional licences and
per-page API pricing put bulk evaluation out of reach for smaller
institutions and independent reviewers, and make cost grow linearly with
volume.

\textbf{P3 -- Latency at institutional scale.} Running a full transformer
classifier on thousands of submissions is a computational bottleneck,
particularly for institutions without GPU hardware. Zero-shot methods such as
DetectGPT are even more expensive per document~\cite{mitchell2023detectgpt}.

\textbf{P4 -- Domain and generator overfitting.} Detectors calibrated on the
output of one generator (for example ChatGPT) lose much of their accuracy on
new generators such as Claude or Llama, on new domains and on paraphrased
text~\cite{dugan2024raid,wang2024m4,li2024mage,krishna2023paraphrasing}.

\textbf{P5 -- Bias and fairness risk.} Perplexity- and diversity-based
heuristics disproportionately flag the writing of non-native English speakers,
whose prose tends to be more predictable to a language model~\cite{liang2023biased}.
This creates ethical and legal risk whenever detector output is used in
academic evaluation.

Fig.~\ref{fig:gapmap} shows how these problems arise from the two existing
families of detectors and how they meet in the deployment gap that this
project addresses.

\begin{figure}[H]
\centering
\begin{tikzpicture}[node distance=0.6cm and 0.6cm]
  \node[blk,text width=4.5cm,minimum height=1.3cm] (heu) {\textbf{Heuristic detectors}\\fast, cheap, interpretable};
  \node[blk,text width=4.5cm,minimum height=1.3cm,right=2.4cm of heu] (tra) {\textbf{Transformer detectors}\\accurate in-domain};
  \node[blk,text width=4.5cm,below=0.7cm of heu] (hl) {Shallow; affected by genre and writer proficiency (P4, P5)};
  \node[blk,text width=4.5cm,below=0.7cm of tra] (tl) {Slow, costly, opaque in commercial form (P1, P2, P3)};
  \node[blk,thick,double,text width=9.0cm,minimum height=1.4cm] (gap) at ($(hl.south)!0.5!(tl.south)+(0,-1.5)$)
      {\textbf{Deployment gap:} institutions must choose between a fast but shallow detector and an accurate but slow, expensive and opaque one};
  \node[blk,text width=9.0cm,below=0.7cm of gap] (vs) {\textbf{VeriScript:} confidence-gated cascade that uses each family where it is most effective};
  \draw[arr] (heu) -- (hl);
  \draw[arr] (tra) -- (tl);
  \draw[arr] (hl.south) -- (gap.north -| hl.south);
  \draw[arr] (tl.south) -- (gap.north -| tl.south);
  \draw[arr] (gap) -- (vs);
\end{tikzpicture}
\caption{Origin of the research gap addressed by VeriScript}
\label{fig:gapmap}
\end{figure}

\subsection{Why a Cascade Is the Appropriate Response}

The difficulty of detection varies widely across documents. A long,
unedited answer produced by a chat assistant with default settings is
usually easy to identify from its perplexity and uniformity alone. A short,
heavily edited or paraphrased text, or a carefully written human essay in a
formal register, is much harder. If most documents in a realistic batch are
easy, then running the most expensive model on all of them wastes most of
the computation.

A cascade addresses P3 directly by spending transformer computation only on
hard documents, and P2 indirectly because an open system that runs on local
hardware has no per-page charge. It addresses P1 because every decision made
on the fast path comes with an explicit feature-level justification. It
mitigates P5 because a heuristic that is uncertain does not decide: borderline
documents, which is where bias in surface statistics is most likely to change
the outcome, are referred to a model that reads the full context. Finally, the
cascade does not solve P4 by itself, but because the thresholds are tuned on
validation data and the evaluation includes held-out generators, the
system's behaviour on unseen sources is measured rather than assumed.

% ---------------------------------------------------------------------
\section{Problem Statement}
% ---------------------------------------------------------------------

\subsection{Informal Statement}

\begin{center}
\fbox{\parbox{0.92\textwidth}{\setstretch{1.3}
\emph{To design, implement and evaluate an open, interpretable system that
classifies an English document as human-written or AI-generated with accuracy
close to that of a fine-tuned transformer classifier, at an average inference
cost close to that of a lightweight heuristic detector, by routing each
document through a confidence-gated cascade, and to verify that the system
generalises to unseen generators and does not disproportionately flag
non-native English writing.}}}
\end{center}

\subsection{Formal Statement}

Let $d$ denote an input document and $y \in \{0,1\}$ its true label, with
$y=0$ for human-written and $y=1$ for AI-generated text. The system must learn
a decision function
\begin{equation}
  f : d \;\mapsto\; \big(\hat{y},\, p(d)\big), \qquad \hat{y}\in\{0,1\},\;\; p(d)\in[0,1],
  \label{eq:fdef}
\end{equation}
where $p(d)$ is a calibrated estimate of $P(y=1 \mid d)$.

The cascade has two component models. The heuristic stage computes a feature
vector $\mathbf{h}(d)$ and a calibrated probability $p_h(d)$ at cost $C_h$;
the transformer stage computes $p_t(d)$ at cost $C_t \gg C_h$. Given an
ambiguity band $[\Tlow,\Thigh]$ with $0 \le \Tlow \le \Thigh \le 1$, the
cascade output is
\begin{equation}
  p(d) \;=\;
  \begin{cases}
     p_h(d), & \text{if } p_h(d) < \Tlow \;\text{ or }\; p_h(d) > \Thigh \quad\text{(fast path)},\\[2pt]
     p_t(d), & \text{if } \Tlow \le p_h(d) \le \Thigh \quad\text{(escalation)},
  \end{cases}
  \label{eq:cascade}
\end{equation}
and $\hat{y} = \mathbb{1}[\,p(d) \ge \tau\,]$ with decision threshold
$\tau$ (0.5 by default). The cost incurred for document $d$ is
\begin{equation}
  C(d) \;=\; C_h + \mathbb{1}\!\left[\Tlow \le p_h(d) \le \Thigh\right] C_t ,
  \label{eq:costd}
\end{equation}
so that the expected cost over the input distribution is
$\mathbb{E}[C] = C_h + \alpha\, C_t$, where
$\alpha = P(\Tlow \le p_h(d) \le \Thigh)$ is the \emph{escalation rate}.

The design problem is to choose the band so as to minimise expected cost
while keeping accuracy close to the transformer-only detector and keeping the
fairness disparity small:
\begin{equation}
\begin{aligned}
  \min_{\Tlow,\,\Thigh}\quad & \mathbb{E}[C(d)] \\
  \text{subject to}\quad & F_1^{\text{cascade}} \;\ge\; F_1^{\text{transformer}} - \varepsilon, \\
  & \Delta_{FPR} \;=\; \big|\,\mathrm{FPR}_{\text{non-native}} - \mathrm{FPR}_{\text{native}}\,\big| \;\le\; \delta ,
\end{aligned}
\label{eq:optproblem}
\end{equation}
where $\varepsilon$ is the largest acceptable loss in F1 relative to the
transformer-only baseline and $\delta$ is the largest acceptable
false-positive-rate disparity. Two limiting cases illustrate the trade-off:
with $\Tlow = 0$ and $\Thigh = 1$ every document is escalated ($\alpha=1$) and
the cascade reduces to the transformer-only detector at cost $C_h + C_t$;
with $\Tlow = \Thigh$ no document is escalated ($\alpha \approx 0$) and the
cascade reduces to the heuristic-only detector at cost $C_h$. The useful
operating points lie between these extremes.

Note that the cost objective is measured relative to the
\emph{transformer-only} detector, which is the system whose accuracy the
cascade aims to retain, and that accuracy is measured relative to the same
baseline; the heuristic-only detector is the lower bound on cost.

\subsection{Assumptions and Constraints}

The formulation rests on the following assumptions.
\begin{enumerate}
  \item Documents are in English and contain at least about 100 words; shorter
        texts carry too little statistical evidence and are reported as
        ``insufficient text''.
  \item A document is treated as entirely human-written or entirely
        AI-generated. Mixed authorship is outside the scope of this phase.
  \item The heuristic stage's probability can be calibrated well enough on a
        validation set for thresholds on it to be meaningful.
  \item Labelled training data covering several generators and domains are
        available from public benchmarks and can be supplemented with custom
        generations.
\end{enumerate}
The system must also satisfy practical constraints: it must run on a
CPU-only machine (with a GPU optional for speed), it must not send documents
to external services, and its components must use openly licensed software and
models.

% ---------------------------------------------------------------------
\section{Depiction of the Problem Statement}
% ---------------------------------------------------------------------

Fig.~\ref{fig:ipo} depicts the problem as a block diagram of inputs,
processes and outputs. Each block is described below.

\begin{figure}[H]
\centering
\begin{tikzpicture}[node distance=0.5cm]
  % Inputs
  \node[blk,text width=3.2cm,align=left,font=\footnotesize] (in1) {Document text\\(.txt, .pdf, .docx, pasted)};
  \node[blk,text width=3.2cm,align=left,font=\footnotesize,below=of in1] (in2) {Thresholds\\$\Tlow$, $\Thigh$, $\tau$};
  \node[blk,text width=3.2cm,align=left,font=\footnotesize,below=of in2] (in3) {Trained models\\(heuristic head, encoder)};
  \node[draw,dashed,thick,fit=(in1)(in2)(in3),inner sep=6pt,
        label={[font=\small\bfseries]above:INPUTS}] (IN) {};
  % Processes
  \node[blk,text width=3.6cm,font=\footnotesize,right=1.4cm of in1] (p1) {1. Pre-processing and tokenisation};
  \node[blk,text width=3.6cm,font=\footnotesize,below=of p1] (p2) {2. Heuristic features and calibrated $p_h(d)$};
  \node[blk,text width=3.6cm,font=\footnotesize,below=of p2] (p3) {3. Confidence gate};
  \node[blk,text width=3.6cm,font=\footnotesize,below=of p3] (p4) {4. Transformer inference $p_t(d)$ (escalated only)};
  \node[blk,text width=3.6cm,font=\footnotesize,below=of p4] (p5) {5. Decision, explanation and logging};
  \node[draw,dashed,thick,fit=(p1)(p5),inner sep=6pt,
        label={[font=\small\bfseries]above:PROCESSES}] (PR) {};
  % Outputs
  \node[blk,text width=3.2cm,align=left,font=\footnotesize,right=1.4cm of p2] (o1) {Label: human / AI};
  \node[blk,text width=3.2cm,align=left,font=\footnotesize,below=of o1] (o2) {Calibrated probability $p(d)$};
  \node[blk,text width=3.2cm,align=left,font=\footnotesize,below=of o2] (o3) {Route taken, latency, feature explanation};
  \node[draw,dashed,thick,fit=(o1)(o2)(o3),inner sep=6pt,
        label={[font=\small\bfseries]above:OUTPUTS}] (OUT) {};
  \draw[arr] (p1) -- (p2);
  \draw[arr] (p2) -- (p3);
  \draw[arr] (p3) -- (p4);
  \draw[arr] (p4) -- (p5);
  \draw[arr] (p3.west) -- ++(-0.35,0) |- (p5.west);
  \draw[arr,very thick] (IN.east |- p2.west) -- (PR.west |- p2.west);
  \draw[arr,very thick] (PR.east |- o2.west) -- (OUT.west |- o2.west);
\end{tikzpicture}
\caption{Inputs, processes and outputs of the problem}
\label{fig:ipo}
\end{figure}

\subsection{Inputs}

The primary input is the text of a document, supplied as a plain-text file,
a PDF or Word file from which text is extracted, or text pasted into the web
interface. A batch input is a folder or a CSV file of documents. The secondary
inputs are configuration values (the ambiguity band $[\Tlow,\Thigh]$ and the
decision threshold $\tau$) and the trained artefacts: the reference language
model used for perplexity, the heuristic classification head with its
calibrator, and the fine-tuned transformer. Table~\ref{tab:inputs} specifies
the inputs.

\begin{table}[H]
\caption{Specification of system inputs}
\label{tab:inputs}
\centering
\small
\begin{tabularx}{\textwidth}{|p{3.4cm}|L|p{3.5cm}|}
\hline
\textbf{Input} & \textbf{Description} & \textbf{Constraint} \\ \hline
Document text & Single document in English & $\ge$ 100 words recommended \\ \hline
Batch file & Folder of files or CSV with \texttt{id}, \texttt{text} columns & UTF-8 encoded \\ \hline
$\Tlow$, $\Thigh$ & Ambiguity band on $p_h(d)$ & $0 \le \Tlow \le \Thigh \le 1$ \\ \hline
$\tau$ & Final decision threshold & Default 0.5 \\ \hline
Reference LM & Small causal LM for perplexity (GPT-2 family) & Loaded locally \\ \hline
Heuristic model & LR / GBDT head with calibrator & Serialised file \\ \hline
Transformer model & Fine-tuned RoBERTa-base or DeBERTa-v3 & Max 512 tokens per window \\ \hline
\end{tabularx}
\end{table}

\subsection{Processes}

\textbf{Pre-processing} normalises Unicode and whitespace, removes markup,
splits the text into sentences and tokens, and rejects inputs that are too
short. \textbf{Heuristic scoring} computes perplexity, burstiness and
stylometric features and passes them through the classification head and the
calibrator to obtain $p_h(d)$. The \textbf{confidence gate} compares $p_h(d)$
with the band: confident documents leave on the fast path, ambiguous
documents are escalated. \textbf{Transformer inference} tokenises escalated
documents, splits long ones into overlapping windows and aggregates the
window probabilities into $p_t(d)$. Finally, the \textbf{decision} step
applies $\tau$, assembles the explanation (the most influential features on
the fast path, or the transformer probability and the heuristic context on the
escalated path) and logs the route and latency for evaluation.

\subsection{Outputs}

For each document the system returns the predicted label, the calibrated
probability $p(d)$, the route taken (fast path or escalated), the time spent in
each stage, and an explanation. In batch mode it also writes a report with
per-document results and summary statistics such as the escalation rate and
throughput. Table~\ref{tab:outputs} specifies the outputs.

\begin{table}[H]
\caption{Specification of system outputs}
\label{tab:outputs}
\centering
\small
\begin{tabularx}{\textwidth}{|p{3.4cm}|L|p{3.5cm}|}
\hline
\textbf{Output} & \textbf{Description} & \textbf{Format} \\ \hline
Label $\hat{y}$ & Human-written or AI-generated & String / integer \\ \hline
Probability $p(d)$ & Calibrated probability of AI authorship & Float in [0,1] \\ \hline
Route & \texttt{fast\_path} or \texttt{escalated} & String \\ \hline
Explanation & Top contributing features with values, or transformer score & JSON object \\ \hline
Latency & Time in heuristic and transformer stages & Milliseconds \\ \hline
Batch report & Per-document results and summary statistics & CSV / JSON \\ \hline
\end{tabularx}
\end{table}

% ---------------------------------------------------------------------
\section{Objectives}
% ---------------------------------------------------------------------

Guided by the research gaps above, the project pursues five objectives.

\begin{enumerate}[label=\textbf{O\arabic*.},leftmargin=1.2cm,itemsep=5pt]
\item \textbf{Heuristic baseline.} Design and implement a statistical detector
      based on perplexity, burstiness and stylometric features (including
      type-token ratio) with a classical machine-learning head (logistic
      regression and gradient-boosted trees), and calibrate its output
      probability.
\item \textbf{Transformer baseline.} Fine-tune a transformer classifier
      (RoBERTa-base, with DeBERTa-v3 as an alternative) on a balanced
      human-versus-AI corpus and establish reference benchmarks for accuracy
      and speed.
\item \textbf{Hybrid architecture.} Develop a confidence-gated cascade that
      returns the heuristic decision for confident documents and escalates
      ambiguous documents to the transformer, with the ambiguity band tuned on
      validation data.
\item \textbf{Unified comparative evaluation.} Measure classification
      performance (Precision, Recall, F1, AUROC) and efficiency (mean and p95
      latency, throughput, model footprint) for the heuristic-only,
      transformer-only and cascade systems, per generator family and in
      aggregate, including a leave-one-generator-out test and a fairness audit
      on non-native English writing.
\item \textbf{System packaging.} Package the final detector as an open-source
      command line tool and web interface suitable for institutional
      evaluation.
\end{enumerate}

Table~\ref{tab:gap-obj} links each problem to the objectives that address it
and to the criterion by which success is judged. The numeric targets are
design goals that were fixed in advance of experimentation and were not
changed afterwards; Chapter~\ref{ch:results} (Table~\ref{tab:criteria-status})
reports which of them the implemented system meets.

\begin{table}[H]
\caption{Mapping of problems to objectives and success criteria}
\label{tab:gap-obj}
\centering
\small
\begin{tabularx}{\textwidth}{|p{3.2cm}|p{1.6cm}|L|}
\hline
\textbf{Problem} & \textbf{Objectives} & \textbf{Success criterion (target)} \\ \hline
P1 Black box & O1, O3, O5 & Every fast-path decision is reported with its top contributing features \\ \hline
P2 Cost and licensing & O5 & Runs locally on open-source software with no per-document charge \\ \hline
P3 Latency at scale & O3, O4 & Mean latency of the cascade at least 50\% below the transformer-only detector on CPU \\ \hline
P4 Generator overfitting & O2, O4 & Cascade F1 on a held-out generator reported and compared with both baselines \\ \hline
P5 Fairness & O3, O4 & False-positive-rate disparity $\Delta_{FPR}$ of the cascade no larger than that of the transformer-only detector \\ \hline
Overall accuracy & O3, O4 & Cascade F1 within $\varepsilon = 0.02$ of the transformer-only detector \\ \hline
\end{tabularx}
\end{table}

% ---------------------------------------------------------------------
\section{Research Questions and Hypotheses}
% ---------------------------------------------------------------------

The objectives are pursued through four research questions, each paired
with a testable hypothesis. The hypotheses were stated before any experiment
was run; Section~\ref{sec:hypotheses} of Chapter~\ref{ch:results} reports
whether the evidence supports or rejects each of them.

\textbf{RQ1 -- Accuracy--cost trade-off.} \emph{How much of the accuracy
advantage of a fine-tuned transformer over a heuristic detector can a
confidence-gated cascade retain, and at what fraction of the transformer's
average inference cost?}\\
\textit{H1:} There exists an ambiguity band for which the cascade's F1 is
within $\varepsilon$ of the transformer-only detector while its escalation rate
$\alpha$ is well below 1, so that its mean latency is substantially lower.

\textbf{RQ2 -- Value of calibration.} \emph{Does calibrating the heuristic
stage improve the cascade's operating points?}\\
\textit{H2:} For the same escalation rate, a cascade built on calibrated
probabilities achieves higher F1 than one built on raw scores, because the
thresholds then select documents on which the heuristic is genuinely
uncertain.

\textbf{RQ3 -- Generalisation.} \emph{How does the cascade behave on a
generator family that was excluded from training and calibration?}\\
\textit{H3:} Accuracy falls for all three systems on the held-out generator,
but the escalation rate of the cascade rises, so that more of the unfamiliar
documents are referred to the stronger model and the cascade degrades no more
than the transformer-only detector.

\textbf{RQ4 -- Fairness.} \emph{Does routing reduce the false-positive
disparity between native and non-native English writing compared with a
heuristic-only detector?}\\
\textit{H4:} Because predictable but human-written essays tend to fall inside
the ambiguity band, the cascade's $\Delta_{FPR}$ is markedly lower than that of
the heuristic-only detector and no higher than that of the transformer-only
detector.

Fig.~\ref{fig:tradeoff} illustrates RQ1 qualitatively. The two baselines are
the end points of the trade-off; the cascade is expected to trace a curve of
operating points between them as the band widens, and the objective is to
find points in the upper-left region, close to the transformer's accuracy at
close to the heuristic's cost.

\begin{figure}[H]
\centering
\begin{tikzpicture}
\begin{axis}[
   width=11cm,height=7.0cm,
   xlabel={Average inference cost per document (relative)},
   ylabel={Detection accuracy (F1, relative)},
   xmin=0,xmax=1.1,ymin=0,ymax=1.1,
   xtick=\empty,ytick=\empty,
   axis lines=left,
   clip=false]
  \fill[pattern=north east lines,pattern color=black!40] (axis cs:0.05,0.82) rectangle (axis cs:0.45,1.02);
  \node[font=\footnotesize,anchor=south] at (axis cs:0.25,1.03) {target region};
  \addplot[only marks,mark=square*,mark size=3pt] coordinates {(0.08,0.45)};
  \node[anchor=west,font=\footnotesize] at (axis cs:0.10,0.45) {heuristic only};
  \addplot[only marks,mark=triangle*,mark size=4pt] coordinates {(1.0,0.95)};
  \node[anchor=south east,font=\footnotesize] at (axis cs:1.02,1.0) {transformer only};
  \addplot[thick,dashed,smooth,mark=*,mark size=1.6pt] coordinates
      {(0.08,0.45) (0.15,0.70) (0.25,0.84) (0.38,0.91) (0.55,0.94) (0.75,0.95) (1.0,0.95)};
  \node[anchor=north west,font=\footnotesize] at (axis cs:0.45,0.86) {cascade as the band widens};
\end{axis}
\end{tikzpicture}
\caption{Qualitative accuracy--cost trade-off motivating the cascade (schematic, not measured data)}
\label{fig:tradeoff}
\end{figure}

% ---------------------------------------------------------------------
\section{Feasibility Study}
% ---------------------------------------------------------------------

\subsection{Technical Feasibility}

All components are built from mature, openly available software and models:
GPT-2 small for perplexity, RoBERTa-base and DeBERTa-v3-base for
classification, scikit-learn and XGBoost for the heuristic head, and PyTorch
with Hugging Face Transformers for neural inference~\cite{paszke2019pytorch,
wolf2020transformers,pedregosa2011sklearn,chen2016xgboost}. Fine-tuning a
base-sized encoder on tens of thousands of documents is routinely done on a
single 16 GB GPU, which is available through free or low-cost cloud notebooks.
The labelled data needed are available from public benchmarks. The project is
therefore technically feasible within the team's resources.

\subsection{Economic Feasibility}

The project uses only open-source software and openly licensed models. The
main costs are GPU time for fine-tuning, which can be met by free cloud tiers,
and API usage for generating custom GPT-4 and Claude texts, which is limited to
a few thousand documents. Llama generations are produced locally. Once built,
the system has no per-document cost, in contrast with commercial detectors
priced per licence or per page.

\subsection{Operational Feasibility}

The intended users interact with the system through a web page or a single
command, and the output is a label, a probability and a short explanation in
plain language. Deployment as a Docker container on an existing institutional
server requires no specialised hardware. Because documents are processed
locally, the system can be used without transferring student work to third
parties, which removes a common obstacle to adoption.

\subsection{Risks and Mitigation}

\begin{table}[H]
\caption{Project risks and mitigation measures}
\label{tab:risks}
\centering
\small
\begin{tabularx}{\textwidth}{|p{4.2cm}|L|}
\hline
\textbf{Risk} & \textbf{Mitigation} \\ \hline
Heuristic stage too weak, so that the escalation rate must be very high & Add stronger single-pass features (e.g.\ conditional probability curvature~\cite{bao2024fast}); report the trade-off curve honestly \\ \hline
Calibration drifts on unseen generators & Evaluate calibration error per generator; recalibrate on a small labelled sample from the new source \\ \hline
Insufficient GPU time for DeBERTa-v3 & RoBERTa-base is the primary model; DeBERTa-v3 is an alternative \\ \hline
Reference LM of the heuristic stage costs as much as the second stage, so that the cascade saves no time & Measure the cost of every stage separately; report the outcome honestly; replace the reference LM by a smaller model if needed \\ \hline
Second-stage model is itself biased or brittle out of domain & Evaluate the transformer on out-of-distribution and fairness sets, not only on the in-distribution test split \\ \hline
API cost of custom generations & Limit generation size; rely on public benchmarks for most training data \\ \hline
Misuse of results as proof of misconduct & Present probabilities with explanations and an explicit advisory note; document limitations \\ \hline
\end{tabularx}
\end{table}

\section{Summary}

This chapter defined the problem addressed by \veriscript{}: building a
detector that is simultaneously accurate, inexpensive, explainable, robust and
fair. The problem was stated formally as the choice of an ambiguity band that
minimises expected inference cost subject to an accuracy constraint relative
to the transformer-only detector and a fairness constraint. The inputs,
processes and outputs were specified, and five objectives with measurable
success criteria were set. The next chapter presents the methodology by which
these objectives will be met.
